\clearpage

\section{Realimentación de estados con acción integral mediante LMI}

De acuerdo con la linealización exacta por realimentación, se definen las
entradas virtuales
\[
  v_x=\ddot{x},\qquad v_z=\ddot{z},\qquad \tau=\ddot{\theta},
\]
Como $m=J=1$. Por tanto, el modelo que se
utiliza para diseñar el controlador es
\[
  \ddot{x}=v_x,\qquad
  \ddot{z}=v_z,\qquad
  \ddot{\theta}=\tau.
\]

Se agrupan las posiciones y velocidades en el vector
\[
  q=
  \begin{bmatrix}x&z&\theta\end{bmatrix}^{T},
  \qquad
  \dot q=
  \begin{bmatrix}\dot{x}&\dot{z}&\dot{\theta}\end{bmatrix}^{T},
\]
y las entradas virtuales en
\[
  v=
  \begin{bmatrix}v_x&v_z&\tau\end{bmatrix}^{T}.
\]
El estado del PVTOL linealizado es
\[
  X=
  \begin{bmatrix}q\\\dot q\end{bmatrix}\in\mathbb{R}^{6}.
\]
Así,
\[
  \dot X=AX+Bv,\qquad y=CX,
\]
donde la salida corresponde a las tres variables que se desean seguir:
\[
  A=
  \begin{bmatrix}
    0_{3\times3}&I_3\\
    0_{3\times3}&0_{3\times3}
  \end{bmatrix},
  \qquad
  B=
  \begin{bmatrix}
    0_{3\times3}\\
    I_3
  \end{bmatrix},
  \qquad
  C=
  \begin{bmatrix}I_3&0_{3\times3}\end{bmatrix}.
\]
De forma explícita,
\[
  A=
  \begin{bmatrix}
  0&0&0&1&0&0\\
  0&0&0&0&1&0\\
  0&0&0&0&0&1\\
  0&0&0&0&0&0\\
  0&0&0&0&0&0\\
  0&0&0&0&0&0
  \end{bmatrix},
  \qquad
  B=
  \begin{bmatrix}
  0&0&0\\
  0&0&0\\
  0&0&0\\
  1&0&0\\
  0&1&0\\
  0&0&1
  \end{bmatrix}.
\]

Para una referencia constante
\[
  r=
  \begin{bmatrix}r_x&r_z&r_\theta\end{bmatrix}^{T},
\]
se define el estado integral
\[
  \eta(t)=\int_0^t(r-y(\tau))\,d\tau
  =
  \begin{bmatrix}\eta_x&\eta_z&\eta_\theta\end{bmatrix}^{T}.
\]
En consecuencia,
\[
  \dot\eta=r-CX.
\]
El sistema aumentado utilizado para el diseño queda:
\[
  X_a=
  \begin{bmatrix}X\\\eta\end{bmatrix}\in\mathbb{R}^{9},
  \qquad
  \dot X_a=A_aX_a+B_av,
\]
con
\[
  \boxed{
  A_a=
  \begin{bmatrix}
    A&0_{6\times3}\\
    -C&0_{3\times3}
  \end{bmatrix}},
  \qquad
  \boxed{
  B_a=
  \begin{bmatrix}B\\0_{3\times3}\end{bmatrix}}.
\]
En forma explícita,
\[
  A_a=
  \begin{bmatrix}
  0&0&0&1&0&0&0&0&0\\
  0&0&0&0&1&0&0&0&0\\
  0&0&0&0&0&1&0&0&0\\
  0&0&0&0&0&0&0&0&0\\
  0&0&0&0&0&0&0&0&0\\
  0&0&0&0&0&0&0&0&0\\
 -1&0&0&0&0&0&0&0&0\\
  0&-1&0&0&0&0&0&0&0\\
  0&0&-1&0&0&0&0&0&0
  \end{bmatrix},
  \qquad
  B_a=
  \begin{bmatrix}
  0&0&0\\
  0&0&0\\
  0&0&0\\
  1&0&0\\
  0&1&0\\
  0&0&1\\
  0&0&0\\
  0&0&0\\
  0&0&0
  \end{bmatrix}.
\]

Se propone la realimentación de estados aumentados
\[
  \boxed{u=-KX_a},
  \qquad K\in\mathbb{R}^{3\times9}.
\]
La matriz $K$ se puede particionar como
\[
  K=\begin{bmatrix}K_q&K_{\dot q}&K_I\end{bmatrix},
\]
con $K_q,K_{\dot q},K_I\in\mathbb{R}^{3\times3}$. En consecuencia,
\[
  u=-K_qq-K_{\dot q}\dot q-K_I\eta.
\]
La dinámica aumentada en lazo cerrado es
\[
  \dot X_a=(A_a-B_aK)X_a.
\]


Se utiliza la función de Lyapunov
\[
  V=X_a^TPX_a,\qquad P=P^T>0.
\]
Para una tasa de decaimiento $\alpha>0$,
\[
  (A_a-B_aK)^TP+P(A_a-B_aK)+\alpha P<0.
\]
Esta desigualdad es bilineal en $P$ y $K$. Se realiza el cambio de variables
\[
  \Gamma=P^{-1},\qquad F=K\Gamma,
\]
de modo que
\[
  K=F\Gamma^{-1}=FP.
\]
La condición equivalente que se resuelve como LMI es
\[
  \boxed{
  A_a\Gamma+\Gamma A_a^T
  -B_aF-F^TB_a^T+\alpha\Gamma<0,
  \qquad \Gamma=\Gamma^T>0.}
\]

